\chapter{Gottes Führung}
\quellseite{8}

Der nachmals so berühmt gewordene Maler Peter v. Cornelius unternahm als Jüngling zusammen mit einem Gefährten eine Fußwanderung durch die Berge. Sehr ermüdet kehrten die beiden am Abend in ein Bauernhaus ein, wo man ihnen gern das erbetene Nachtquartier gewährte. Doch ehe die Hausgenossen sich nach einem kräftigen Abendbrot zur Ruhe begaben, versammelte der Bauer noch die Familie und das Gesinde zur gewohnten Abendandacht; auch die jungen Gäste nahmen, um die freundlichen Wirthe nicht zu kränken, daran Theil, obgleich man ihnen wohl anmerken konnte, daß sie nicht viel danach fragten. Als man sich dann trennte, war es, trotz der Anstrengungen des Tages, den Malern noch zu früh, sich schon schlafen zu legen, und Cornelius trat hinaus auf die hölzerne Gallerie, die nach der dortigen Bauart das Haus umgab. Er genoß dankbar die herrliche, erfrischende Luft und bewunderte den tiefdunkeln Nachthimmel, an dem strahlend hell zahllose Sterne funkelten. Wie er so in Gedanken versunken dastand, hörte er in einem der auf die Gallerie mündenden Zimmer sprechen. Unwillkürlich horchte der junge Maler auf. Das war die Stimme des alten Bauern, der beim Vater im Himmel noch einmal sich und die Seinen zu treuer Obhut empfahl. Zum Schluß gedachte er auch fürbittend der unter seinem gastlichen Dache eingekehrten Fremden, die wohl noch recht wenig vom lieben Gott und seinem Wort zu wissen schienen; der Herr möchte sich doch auch an ihnen nicht unbezeugt lassen! Wie gingen dem unfreiwilligen Zuhörer diese Worte des schlichten, frommen Peters ans Herz! War es nicht schon ein Beweis von Gottes suchender Liebe, daß er ihn und seinen Gefährten gerade in dies Haus geführt und zum Zeugen dieses Gebets gemacht hatte? — Von diesem Abend an verwandelte sich bei dem jungen Cornelius die frühere Gleichgültigkeit gegen religiöse Dinge in ein Suchen nach rechter Erkenntniß; und Gott ließ es dem Aufrichtigen gelingen. Die im schlichten Bauernhaus empfangenen Eindrücke wurden von bestimmendem Einfluß auf das Leben des späteren großen Künstlers.
